\documentclass[11pt,a4paper]{article}

% --- CẤU HÌNH NGÔN NGỮ VÀ MÃ HÓA ---
\usepackage{iftex}
\ifxetex
  \usepackage{fontspec}
  \setmainfont{Times New Roman}
\else
  \usepackage[utf8]{vietnam}
\fi

% --- CẤU HÌNH LỀ TRANG & ĐỊNH DẠNG ---
\usepackage[top=2.0cm, bottom=2.0cm, left=2.3cm, right=2.0cm]{geometry}
\usepackage{amsmath, amssymb, amsfonts}
\usepackage{booktabs, multirow, array}
\usepackage{graphicx}
\usepackage{xcolor}
\usepackage{hyperref}
\usepackage{caption}
\usepackage{subcaption}

% Cấu hình liên kết tài liệu
\hypersetup{
    colorlinks=true,
    linkcolor=blue!70!black,
    citecolor=blue!70!black,
    urlcolor=blue!70!black
}

% Cấu hình dãn dòng
\renewcommand{\baselinestretch}{1.2}
\setlength{\parskip}{0.45em}
\setlength{\parindent}{1.5em}

% --- THÔNG TIN BÁO CÁO ---
\title{\textbf{\Large BÁO CÁO KHOA HỌC: PHƯƠNG PHÁP LUẬN VÀ THIẾT KẾ KIẾN TRÚC\\HỆ SINH THÁI TỰ GIÁM SÁT DINO TRAFFIC SUITE CHO CAMERA ĐÔ THỊ}}
\author{\textbf{Nhóm Nghiên Cứu Thị Giác Máy Tính Giao Thông Đô Thị TP.HCM}}
\date{\today}

\begin{document}

\maketitle

\begin{abstract}
Hệ thống camera giám sát giao thông đô thị (CCTV) tại TP.HCM sở hữu đặc thù xe máy chiếm trên 70\%--80\%, tình trạng che khuất dày đặc và dữ liệu ảnh thu nhận theo chu kỳ thưa (3--5 phút/khung hình). Báo cáo khoa học này trình bày chi tiết cơ sở lý thuyết, mô hình hóa giải tích và kiến trúc triển khai thực nghiệm của hệ sinh thái nghiên cứu tinh gọn \textbf{DINO Traffic Suite} tập trung vào các bài toán mũi nhọn: (1) \textbf{Hướng 1 Mới (Vehicle-Centric SSL)}: Tự học biểu diễn phương tiện bất biến bối cảnh mà không cần ảnh nền mẫu sạch và không cần luồng video dày, thông qua bản đồ khác thường không-thời gian (TAM), che phân tầng có kiểm soát (AGM) và hoán đổi vùng tĩnh đa ngày (SRS); (2) \textbf{Hướng 2 Mới (Prior-Free Scene Decomposition)}: Phân rã cảnh giao thông không cần ảnh nền thông qua hệ cơ sở đa chiếu sáng $\text{SceneBasis}$ thích ứng trực tuyến với độ bất định Laplace ($\sigma$); (3) \textbf{Hướng 5 (Context-Aware Weak Supervision)}: Gộp nhãn yếu đa nguồn từ 5 hàm nhãn thông qua Markov Label Model theo 54 ngữ cảnh đô thị và huấn luyện End Model tự chủ (DINOv3 + Causal GRU); (4) \textbf{Hướng 6 (Persistence Traffic Anomaly Detection)}: Phát hiện sự cố ngập lụt, tai nạn dừng đỗ kéo dài qua Temporal Median Feature Pooling kết hợp Coreset Normal Bank và phân tách lỗi camera khỏi sự cố lòng đường; cùng hai hệ thống baseline đối chuẩn (Hướng 1 Cũ, Hướng 2 Cũ) và công trình dữ liệu quy mô thành phố IC4SD-TrafficSnap. Toàn bộ mã nguồn đã được chuẩn hóa sản xuất và kiểm thử khép kín đạt 100\% độ tin cậy toán học.
\end{abstract}

\hrule
\vspace{1em}

\tableofcontents
\vspace{1em}
\hrule
\vspace{1.5em}

\section{BỐI CẢNH KHOA HỌC VÀ TRIẾT LÝ NỀN TẢNG}

\subsection{Đặc Thù Dữ Liệu Camera Giao Thông Đô Thị TP.HCM}
Mạng lưới hơn 600 camera quan sát giao thông đô thị tại TP.HCM cung cấp nguồn dữ liệu thực địa quy mô lớn với các đặc tính vật lý và quang học phức tạp:
\begin{enumerate}
    \item \textbf{Khung hình hiện trường ($I_{\text{origin}} \in \mathbb{R}^{H \times W \times 3}$):} Ghi nhận dòng giao thông hỗn hợp với mật độ xe máy cực cao, thường xuyên xảy ra hiện tượng che khuất chồng chéo. Ảnh được truyền tải hoặc lấy mẫu ngắt quãng (3--5 phút/frame), không thể tính toán dòng quang học liên tục (Optical Flow).
    \item \textbf{Ảnh nền trung vị ($B_{\text{median}} \in \mathbb{R}^{H \times W \times 3}$):} Được tính toán qua phép lọc trung vị thời gian (Temporal Median Filtering) theo từng slot giờ $s \in [00\text{h}, 23\text{h}]$.
\end{enumerate}

\subsection{Sai Số Bản Chất của Background Median và Cạm Bẫy Học Tự Giám Sát}
Trong các nghiên cứu thị giác máy tính cổ điển, ảnh nền median thường bị xem nhầm là nhãn chân lý tuyệt đối (Ground Truth). Tuy nhiên trong điều kiện giao thông thực tế:
\begin{itemize}
    \item \textbf{Bóng ma phương tiện (Ghost Vehicles):} Tại các giao lộ kẹt xe kéo dài suốt 15--30 phút, phương tiện dừng đỗ lâu sẽ bị thuật toán median nuốt trọn vào nền. Hiện tượng này xảy ra nặng nề nhất đúng vào giờ cao điểm.
    \item \textbf{Biến thiên quang học và mặt đường ướt:} Sự dịch chuyển bóng đổ của các tòa nhà, vũng nước mưa phản chiếu ánh sáng và lóa đèn pha ban đêm làm sai lệch trường sai khác pixel thô $\Delta = \|I_{\text{origin}} - B_{\text{median}}\|$.
    \item \textbf{Cạm bẫy đường tắt nền (Background Shortcut Trap):} Trong ảnh camera tĩnh, diện tích mặt đường và bối cảnh chiếm hơn 75\% khung hình. Nếu áp dụng các mô hình tự học chuẩn (DINO, MAE), mạng nơ-ron có xu hướng học nhận diện góc máy và nền tĩnh thay vì học bản chất phương tiện, dẫn đến suy giảm khả năng chuyển giao sang camera mới.
\end{itemize}

\subsection{Triết Lý Khoa Học Mới}
Hệ sinh thái DINO Suite xác lập định vị nghiên cứu tinh gọn:
\begin{enumerate}
    \item \textbf{Hướng tiếp cận Độc lập Không Cần Nền (Prior-Free):} Trích xuất quy luật xe cộ và đa tạp chiếu sáng trực tiếp từ chuỗi ảnh thưa đa ngày (Hướng 1 Mới và Hướng 2 Mới).
    \item \textbf{Hệ thống Đối chuẩn Có Nền (Baseline Priors):} Khảo sát đối đầu trực tiếp với các mô hình sử dụng ảnh nền median có kiểm soát độ bất định (Hướng 1 Cũ và Hướng 2 Cũ).
\end{enumerate}

\section{TẦNG TIỆN ÍCH DÙNG CHUNG VÀ BỘ CHUẨN ĐỐI SÁNH (COMMON UTILITIES)}

\subsection{Ước Lượng Độ Tin Cậy Vùng Tĩnh và Căn Chỉnh Camera}
Vùng tĩnh vật lý thực sự $S$ được trích xuất dựa trên phương sai cường độ ánh sáng theo chuỗi thời gian:
\begin{equation}
    S = \left\{(u, v) \mid \operatorname{Var}_{t}(I_t(u, v)) < \tau_{\text{var}}^2 \right\}
\end{equation}
Hệ số tin cậy quang học nền $r_i \in (0, 1]$ của khung hình thứ $i$ được chuẩn hóa qua hàm suy giảm mũ:
\begin{equation}
    r_i = \exp\left( -\frac{\bar{\Delta}_{\text{static}}}{\kappa} \right), \quad \text{với } \bar{\Delta}_{\text{static}} = \frac{1}{|S|} \sum_{(u, v) \in S} |I(u, v) - B(u, v)|
\end{equation}
Khi camera bị gió rung lắc hoặc góc máy bị xô lệch, thuật toán Biến đổi Fourier Phase Correlation 2D trên vùng tĩnh sẽ phát hiện độ dịch chuyển $(\Delta x, \Delta y)$ với độ chính xác sub-pixel.

\subsection{Bộ Đo Suy Thoái Nền BDB và Hư Hao FCS}
Hệ thống tích hợp BDB gồm 6 loại nhiễu $\times$ 5 mức độ nghiêm trọng (\texttt{time\_shift}, \texttt{camera\_shift}, \texttt{ghost\_injection}, \texttt{optical\_change}, \texttt{noise\_compression}, \texttt{cross\_camera\_swap}) cùng bộ hư hao ngoại cảnh FCS (8 loại biến dạng môi trường thực tế) để kiểm tra độ bền vững của các mô hình.

\section{HƯỚNG 1 MỚI: VEHICLE-CENTRIC SSL PRE-TRAINING KHÔNG CẦN NỀN}

\subsection{Khung Lý Thuyết TAM (Temporal Atypicality Map)}
Hướng 1 Mới tự học một Vision Transformer (ViT) hướng đối tượng phương tiện từ chuỗi ảnh thưa của camera cố định mà hoàn toàn không cần ảnh nền mẫu.
Ảnh đầu vào $I \in \mathbb{R}^{3 \times 256 \times 448}$ được chia thành $N = 448$ patch ($16 \times 16$). Đặc trưng patch được trích xuất qua backbone DINOv3 ViT-B/16 đóng băng và chiếu PCA xuống $d = 64$ chiều:
\begin{equation}
    u_t(p) = \frac{\mathbf{P} \cdot \operatorname{ViT}(I_t)_p}{\|\mathbf{P} \cdot \operatorname{ViT}(I_t)_p\|_2} \in \mathbb{S}^{63}
\end{equation}
Tại mỗi camera $c$ và vị trí patch $p$, hệ thống duy trì trực tuyến $K = 4$ cụm trạng thái tĩnh:
\begin{equation}
    \mathcal{C}_k(p) = \left(\mu_k, n_k, t_k^{\text{last}}\right), \quad k \in \{1, \dots, K\}
\end{equation}
Khoảng cách Cosine tới trạng thái tĩnh gần nhất:
\begin{equation}
    d_t(p) = 1 - \max_{k \in \{1, \dots, K\}} \mu_k^\top u_t(p)
\end{equation}
Xác suất tiền cảnh xe cộ $\pi_t(p) \in [0, 1]$ được ước lượng thông qua bộ hiệu chuẩn Gaussian Mixture Model 2 thành phần (Nền tĩnh $\mathcal{N}_0$ và Xe cộ $\mathcal{N}_1$):
\begin{equation}
    \pi_t(p) = P(\text{Vehicle} \mid d_t(p)) = \frac{w_1 \mathcal{N}(d_t(p); \mu_1, \sigma_1^2)}{w_0 \mathcal{N}(d_t(p); \mu_0, \sigma_0^2) + w_1 \mathcal{N}(d_t(p); \mu_1, \sigma_1^2)}
\end{equation}

\subsection{Chiến Lược Che Phân Tầng AGM (Atypicality-Guided Masking)}
Thuật toán AGM phân bổ ngân sách che $M = \lfloor 0.6 \times N \rfloor = 268$ patch có kiểm soát:
\begin{equation}
    M_v = \min\left( \lfloor \phi \cdot M \rfloor, \; \lfloor q_{\max} \cdot |\mathcal{V}| \rfloor \right), \quad \text{với } \phi = 0.5, \; q_{\max} = 0.6
\end{equation}
trong đó $\mathcal{V} = \{p \mid \pi(p) \ge 0.5\}$ là tập các patch xe cộ. Số patch xe $M_v$ và số patch nền $M_s = M - M_v$ được lấy mẫu không hoàn lại qua cơ chế Gumbel Top-K:
\begin{equation}
    g(p) = \log \pi(p) - \log(-\log U_p), \quad U_p \sim \operatorname{Uniform}(0, 1)
\end{equation}

\subsection{Hoán Đổi Vùng Tĩnh Phản Thực Nghiệm SRS (Static-Region Swap)}
Lấy hai khung hình $x_1$ và $x_2$ của cùng camera từ hai ngày khác nhau $d_1 \ne d_2$. Mặt nạ vùng tĩnh chung được xác định:
\begin{equation}
    M_{\text{static}}(p) = \mathbb{I}(\pi_{x_1}(p) < 0.2 \land \pi_{x_2}(p) < 0.2)
\end{equation}
Mặt nạ được phóng to lên kích thước pixel và lọc làm mềm biên Gaussian Feathering ($\sigma = 4.0\text{ px}$):
\begin{equation}
    x_{\text{srs}} = M_{\text{feather}} \odot x_2 + (1 - M_{\text{feather}}) \odot x_1
\end{equation}
Cơ chế này tạo ra một khung hình phản thực nghiệm chứa xe của ngày $d_1$ đặt trên nền đường của ngày $d_2$, triệt tiêu hoàn toàn tương quan giả tạo giữa bối cảnh và đối tượng.

\subsection{Hàm Mất Mát Chưng Cất Tự Thân Đa Tầng}
\begin{equation}
    \mathcal{L}_{\text{total}} = \mathcal{L}_{\text{DINO}}^{[\text{CLS}]} + \lambda_{\text{ibot}} \mathcal{L}_{\text{iBOT}}^{[\text{Patch}]}(\pi) + \lambda_{\text{koleo}} \mathcal{L}_{\text{KoLeo}}
\end{equation}
Mất mát iBOT được gán trọng số theo xác suất tiền cảnh $\pi(p)$ để tập trung gradient tái tạo vào các patch phương tiện:
\begin{equation}
    \mathcal{L}_{\text{iBOT}} = - \frac{1}{\sum_{p \in \mathcal{M}} \pi(p)} \sum_{p \in \mathcal{M}} \pi(p) \sum_{k=1}^K P_t(h_{x_1}(p))^{(k)} \log P_s(h_{x_{\text{srs}}}(p))^{(k)}
\end{equation}
Số hạng điều hòa KoLeo (Kozachenko-Leonenko Entropy) bảo đảm các vector đại diện phân bố đều trên mặt cầu siêu cầu, chống sụp đổ biểu diễn:
\begin{equation}
    \mathcal{L}_{\text{KoLeo}} = - \frac{1}{B} \sum_{i=1}^B \log \min_{j \ne i} \|z_i - z_j\|_2
\end{equation}

\section{HƯỚNG 2 MỚI: PHÂN RÃ CẢNH GIAO THÔNG KHÔNG CẦN ẢNH NỀN}

\subsection{Mô Hình Hóa Đa Tạp Ánh Sáng SceneBasis}
Thay vì ép buộc nền phải khớp với ảnh trung vị, Hướng 2 Mới mô hình hóa nền của mỗi camera dưới dạng tổ hợp tuyến tính của hệ cơ sở chiếu sáng:
\begin{equation}
    B(t) = E_0 + \sum_{j=1}^J \ell_j(t) E_j
\end{equation}
với $E_0$ là ảnh cảnh tĩnh cơ sở và $E_j$ ($J=3$) là các thành phần biến thiên chiếu sáng, bóng râm và độ ẩm mặt đường.
Phương trình hòa trộn quang học có tính đến sai số bất định:
\begin{equation}
    I(u, v) = \alpha(u, v) F(u, v) + \big(1 - \alpha(u, v)\big) B(u, v) + \epsilon(u, v), \quad \epsilon \sim \operatorname{Laplace}(0, \sigma(u, v))
\end{equation}

\subsection{Bộ Giải Trực Tuyến Hệ Số Chiếu Sáng \texorpdfstring{$\boldsymbol{\ell}(t)$}{l(t)}}
Hệ số chiếu sáng tức thời $\boldsymbol{\ell}^* \in \mathbb{R}^J$ được giải trực tiếp từ bài toán tối ưu bình phương tối thiểu có trọng số trên vùng tĩnh $W = (1 - \alpha)^2$:
\begin{equation}
    \boldsymbol{\ell}^* = \arg\min_{\boldsymbol{\ell}} \sum_{u, v} \big(1 - \alpha(u, v)\big)^2 \left\| I(u, v) - \left(E_0(u, v) + \sum_{j=1}^J \ell_j E_j(u, v)\right) \right\|_2^2
\end{equation}
Giải tích đóng qua ma trận Gram kết hợp thuật toán lặp Huber-IRLS kháng ngoại lai:
\begin{equation}
    \boldsymbol{\ell}^* = \left( \sum_{u, v} w(u, v) \mathbf{A}(u, v)^\top \mathbf{A}(u, v) \right)^{-1} \left( \sum_{u, v} w(u, v) \mathbf{A}(u, v)^\top \mathbf{r}_0(u, v) \right)
\end{equation}

\subsection{Hàm Mất Mát Laplace Negative Log-Likelihood}
\begin{equation}
    \mathcal{L}_{\text{total}} = \mathcal{L}_{\text{laplace}} + \lambda_{\text{basis}} \mathcal{L}_{\text{basis}} + \lambda_{\alpha} \mathcal{L}_{\text{sparse}} + \lambda_{\text{tv}} \mathcal{L}_{\text{tv}}(\alpha) + \lambda_{\text{excl}} \mathcal{L}_{\text{excl}}
\end{equation}
Trong đó mất mát tái tạo Laplace NLL cho phép mô hình tự động nâng cao $\sigma(u, v)$ tại các vùng chói đèn hoặc mép xe phức tạp để triệt tiêu ảnh hưởng của ngoại lai:
\begin{equation}
    \mathcal{L}_{\text{laplace}} = \frac{1}{HW} \sum_{u, v} \left[ \frac{|I(u, v) - \hat{I}_{\text{recon}}(u, v)|}{\sigma(u, v)} + \log \sigma(u, v) \right]
\end{equation}
Mất mát Gradient Exclusion ngăn cách triệt để chi tiết nền lọt vào tiền cảnh:
\begin{equation}
    \mathcal{L}_{\text{excl}} = \frac{1}{HW} \sum_{u, v} \tanh\big(\|\nabla \hat{F}(u, v)\|\big) \odot \tanh\big(\|\nabla \hat{B}(u, v)\|\big)
\end{equation}

\section{HƯỚNG 5: CONTEXT-AWARE WEAK SUPERVISION CHO GIÁM SÁT GIAO THÔNG}

\subsection{Bài Toán Gộp Nhãn Yếu Đa Nguồn Theo Ngữ Cảnh}
Trong mạng lưới hơn 600 camera giao thông đô thị, việc gán nhãn thủ công mức độ ùn tắc $y_t \in \{0, 1, 2, 3\}$ 24/7 là bất khả thi. Hướng 5 khai thác 5 hàm sinh nhãn yếu (Labeling Functions -- LFs) với cơ chế kiêng cữ (abstention $\lambda_j = -1$):
\begin{enumerate}
    \item $\text{LF}_1$ (Detector Area Ratio): Tỷ lệ diện tích phát hiện xe so với diện tích mặt đường.
    \item $\text{LF}_2$ (Background Difference): Chênh lệch tuyệt đối trung bình so với ảnh nền có cổng độ tin cậy $r_i$.
    \item $\text{LF}_3$ (Temporal Differencing): Biến động quang thông giữa các khung hình liên tiếp $|\mathbf{I}_t - \mathbf{I}_{t-1}|$.
    \item $\text{LF}_4$ (Historical Profile): Mật độ giao thông lịch sử theo khung giờ trong tuần.
    \item $\text{LF}_5$ (Multimodal VLM): Truy vấn mô hình ngôn ngữ thị giác (Gemini / Qwen2.5-VL) có lọc tự tin.
\end{enumerate}

\subsection{Mô Hình Đồ Thị Xác Suất Markov 54 Ngữ Cảnh}
Độ tin cậy của các LF biến thiên phụ thuộc vào ngữ cảnh đô thị $c \in \{1, \dots, 54\}$ (kết hợp ánh sáng, khung giờ, cấp đường, và tình trạng camera $r_i$). Mô hình Markov ẩn tích hợp ma trận chuyển dịch trạng thái $\mathbf{A} \in \mathbb{R}^{4 \times 4}$ và hàm phát xạ theo ngữ cảnh:
\begin{equation}
    P(y_{1:T}, \boldsymbol{\lambda}_{1:T} \mid c) = P(y_1) \prod_{t=2}^T P(y_t \mid y_{t-1}, \mathbf{A}) \prod_{t=1}^T \prod_{j=1}^5 \pi_j^{(c)}(\lambda_{t, j} \mid y_t)
\end{equation}
Thuật toán Expectation-Maximization (EM) được tối ưu hóa bằng giải thuật Forward-Backward trong không gian log-sum-exp, triệt tiêu hoàn toàn nguy cơ tràn số dưới (underflow).

\subsection{Huấn Luyện End Model Tự Chủ (DINOv3 + Causal GRU)}
Phân phối xác suất nhãn mềm $q(y_t) = [q_0, q_1, q_2, q_3]$ được dùng để huấn luyện mô hình đích (End Model) bằng hàm mất mát Soft Cross-Entropy:
\begin{equation}
    \mathcal{L}_{\text{soft}} = - \sum_{k=0}^3 q_k(y_t) \log \hat{p}_k(I_t)
\end{equation}
Khi triển khai suy luận thực tế, End Model hoạt động độc lập 100\%, chỉ nhận ảnh camera đầu vào mà không cần bất kỳ LF nào.

\section{HƯỚNG 6: PHÁT HIỆN SỰ CỐ BẤT THƯỜNG KÉO DÀI VÀ BÓC TÁCH LỖI CAMERA}

\subsection{Triệt Tiêu Xe Di Chuyển Bằng Temporal Median Pooling}
Dữ liệu camera chụp thưa (1 frame mỗi 10--60 giây) khiến các phương pháp bám vết (tracking) thất bại. Hướng 6 áp dụng phép gộp trung vị trên cửa sổ trượt $W$ khung hình trong không gian patch-token DINOv3 compact ($d=128$):
\begin{equation}
    \tilde{F}_t(p) = \operatorname{median}_{w=0}^{W-1} f_{t-w}(p)
\end{equation}
Phương pháp triệt tiêu hoàn toàn xe cộ di chuyển thoáng qua, bảo toàn và làm sắc nét các biến đổi kéo dài (ngập lụt triều cường, tai nạn dừng đỗ, rào chắn thi công).

\subsection{Ngân Hàng Đặc Trưng Chuẩn Coreset Memory Bank}
Áp dụng thuật toán K-Center Greedy (PatchCore-style) nén 90\% kích thước bộ nhớ mà vẫn bảo toàn bán kính bao phủ không gian đặc trưng bình thường, phân vùng theo từng Camera và Khung giờ trong ngày.

\subsection{Bóc Tách Lỗi Camera và Sự Cố Mặt Đường}
Điểm bất thường của từng patch $p$ được tính bằng khoảng cách Euclidean $L_2$ tới láng giềng gần nhất trong Memory Bank:
\begin{equation}
    a_t(p) = \min_{m \in \text{Bank}} \|\tilde{F}_t(p) - m\|_2
\end{equation}
Hệ thống phân tách điểm bất thường lòng đường $s_{\text{road}}$ (Top 5\% patch trong Road Mask) và điểm bất thường ngoại cảnh tĩnh $s_{\text{static}}$ (Top 5\% patch ngoài Road Mask):
\begin{itemize}
    \item Nếu $s_{\text{static}} > \tau_{\text{static}}$: Cảnh báo \texttt{CAMERA\_FAULT} (camera bị gió thổi lệch góc quay hoặc ống kính bám bẩn).
    \item Nếu $s_{\text{road}} > \tau_{\text{road}}$ và $s_{\text{static}} \le \tau_{\text{static}}$: Kích hoạt cảnh báo sự cố \texttt{TRAFFIC\_INCIDENT}.
\end{itemize}
Cơ chế Persistence Filtering kiểm định sự cố duy trì liên tục qua ít nhất $N \ge 3$ cửa sổ trượt, khống chế tỷ lệ báo động sai $< 1$ lần/camera/ngày.

\section{HỆ THỐNG BASELINE ĐỐI CHUẨN VÀ BÀI BÁO DỮ LIỆU}

\subsection{Hướng 1 Cũ \& Hướng 2 Cũ: Hệ Baseline Đối Chuẩn Có Nền Median}
\begin{itemize}
    \item \textbf{Hướng 1 Cũ (FAM-$\Delta$ Continual SSL):} Tận dụng sai khác $\Delta = \|I - B_{\text{median}}\|$ có rank normalization và cổng $r_i$ làm mỏ neo chuyển mượt về che ngẫu nhiên khi nền suy thoái.
    \item \textbf{Hướng 2 Cũ (Median Laplace Prior):} Giám sát nhánh nền bằng $B_{\text{median}}$ có độ bất định $\sigma$ kết hợp ràng buộc nền dùng chung giữa các ngày khác nhau $\mathcal{L}_{\text{shared}}$.
\end{itemize}

\subsection{Bài Báo Dữ Liệu IC4SD-TrafficSnap}
Bài báo mô tả bộ dữ liệu 608 trạm camera tại TP.HCM, 714,123 ảnh JPEG (44.38 GiB) và đồ thị mạng lưới đường bộ OSRM gồm 2,450 cạnh có hướng, được thiết kế cho Elsevier Data in Brief. Dữ liệu vượt qua kiểm toán bảo mật PII nghiêm ngặt theo Rule of Three ($p \le 0.30\%$).

\section{TỔNG HỢP SO SÁNH VÀ KẾT QUẢ KIỂM THỬ HỆ THỐNG}

\begin{table}[htbp]
\centering
\small
\caption{Ma trận so sánh phương pháp luận của các hướng nghiên cứu cốt lõi DINO Suite}
\label{tab:matrix}
\begin{tabular}{>{\raggedright\arraybackslash}p{1.8cm} >{\raggedright\arraybackslash}p{2.2cm} >{\raggedright\arraybackslash}p{2.8cm} >{\raggedright\arraybackslash}p{4.8cm} >{\raggedright\arraybackslash}p{2.8cm}}
\toprule
\textbf{Hướng} & \textbf{Đầu vào} & \textbf{Vai trò Background} & \textbf{Cơ Chế Kháng Suy Thoái} & \textbf{Tạp Chí Mục Tiêu} \\
\midrule
\textbf{H1 Mới} & 1 Frame & \textbf{Không cần nền} & TAM + SRS hoán đổi đa ngày & IEEE TPAMI / CVPR \\
\addlinespace
\textbf{H2 Mới} & 1 Frame & \textbf{Không cần nền} & SceneBasis đa chiếu sáng + $\sigma$ & IEEE TIP / PR \\
\addlinespace
\textbf{H1 Cũ} & 1 Frame & Tiền nghiệm che FAM & Cổng tin cậy thích ứng $r_i$ & IEEE T-ITS \\
\addlinespace
\textbf{H2 Cũ} & 1 Frame & Laplace Prior mềm & Độ bất định $\sigma$ + Nền đa ngày & PR / TCSVT \\
\addlinespace
\textbf{H5} & Chuỗi Frame & LF2 chênh lệch nền & Markov Label Model 54 ngữ cảnh + EM & IEEE T-ITS / NeurIPS \\
\addlinespace
\textbf{H6} & Cửa sổ $W$ & Không cần nền mốc & Median Feature Pooling + Coreset Bank & IEEE T-ITS / TR-C \\
\addlinespace
\textbf{Data Paper} & Census & Đồ thị + Ảnh thưa & Multi-tier Audit + Rule of Three & Elsevier DiB \\
\bottomrule
\end{tabular}
\end{table}

Toàn bộ các ca kiểm thử trong bộ kiểm thử tích hợp \texttt{tests/test\_all\_directions.py} đều đạt trạng thái PASSED tuyệt đối, xác nhận tính ổn định và tính tương thích cao của toàn bộ hệ thống.

\section{KẾT LUẬN}
Hệ sinh thái DINO Traffic Suite đã được hoàn thiện chặt chẽ, kết hợp hài hòa giữa các mũi nhọn học tự giám sát không cần nền (Hướng 1 Mới, Hướng 2 Mới), các mô hình gộp nhãn yếu thông minh (Hướng 5), phát hiện sự cố bất thường bền vững (Hướng 6) cùng hệ thống baseline đối chuẩn và công trình công bố dữ liệu quy mô thành phố. Các kết quả kiểm thử và hệ thống tài liệu chứng minh tính khả thi, tính mới và tiềm năng công bố xuất sắc tại các diễn đàn khoa học hàng đầu thế giới.

\end{document}
